\documentclass[]{standalone}
\usepackage[T1]{fontenc}\usepackage{tikz}
\usepackage{amsmath, amsfonts}
\usetikzlibrary{arrows, shapes.geometric, calc, positioning, decorations.pathreplacing}

% overwrite (either \providecommand or \renewcommand) to control the parts of the diagram which are drawn or not.
\providecommand\DrawBtcState{1}
\providecommand\DrawEthState{1}
\providecommand\DrawBtcRefl{1}
\providecommand\DrawEthRefl{1}
\providecommand\DrawPoRLabel{1}
\providecommand\DrawPoRInState{1}
\providecommand\DrawEthLT{3}
\providecommand\DrawBtcLT{4}

\begin{document}
% default params for image
\providecommand\DrawBtcState{0}%
\providecommand\DrawEthState{0}%
\providecommand\DrawBtcRefl{0}%
\providecommand\DrawEthRefl{0}%
\providecommand\DrawPoRLabel{0}%
\providecommand\DrawPoRInState{0}%
\providecommand\DrawEthLT{3}%
\providecommand\DrawBtcLT{4}%
\providecommand\NameChainA{Bitcoin}
\providecommand\AbbrChainA{BTC}
\providecommand\NameChainB{Ethereum}
\providecommand\AbbrChainB{ETH}
\providecommand\ChainBBlocksN{40}
%
% define \ifnodedefined
\makeatletter % https://tex.stackexchange.com/questions/37709/how-can-i-know-if-a-node-is-already-defined
\long\def\ifnodedefined#1#2#3{%
    \@ifundefined{pgf@sh@ns@#1}{#3}{#2}%
}%
\makeatother%
\tikzstyle{chainName} = [rectangle, rounded corners, minimum width=3cm, minimum height=1cm, text centered, draw=black]
\tikzstyle{block} = [rectangle, rounded corners, minimum width=1cm, minimum height=1cm, draw=black]
\tikzstyle{dots} = [inner sep=0pt, text height=10pt]
\tikzstyle{arrowParent} = [thick, ->, -to, shorten >= 2pt, shorten <= 2pt]
\tikzstyle{arrowRefl} = [dashed, ->, -to, shorten >= 2pt, shorten <= 2pt]
\tikzstyle{braceLL} = [decorate, decoration={brace,amplitude=6pt,mirror,raise=2pt}]
\tikzstyle{braceLR} = [decorate, decoration={brace,amplitude=6pt,raise=-8pt}]
\tikzstyle{ethBraceL} = [braceLL, sloped, rotate=180]
\tikzstyle{braceRR} = [decorate, decoration={brace,amplitude=6pt,raise=2pt}]
\tikzstyle{braceRL} = [decorate, decoration={brace,amplitude=6pt,raise=-8pt,mirror}]
\tikzstyle{ethBraceR} = [braceRR, sloped]
%
\begin{tikzpicture}[node distance=0.8cm]
    \pgfmathtruncatemacro\lastBtc{round(\DrawBtcLT-1)}
    \pgfmathtruncatemacro\lastEth{round(\DrawEthLT-1)}

    \def\ethE{ethS}
    \def\ethAll{ethS}
    \def\ethPorHeaders{header}
    \def\ethPorIncl{PoR}
    \ifnum 1<\ChainBBlocksN
        \def\ethE{ethE}
        \def\ethAll{ethAll}
        \def\ethPorHeaders{\ChainBBlocksN\:headers}
        \def\ethPorIncl{\ChainBBlocksN\:PoRs}
    \fi

    \node (btcChain) [chainName] {\NameChainA};
    \node (ethChain) [chainName, right of=btcChain, xshift=4cm] {\NameChainB};

    \foreach \net/\oNet in {btc/eth,eth/btc} {
        \node (\net Dots) [below of=\net Chain, yshift=0cm] {$\vdots$};
    }

    \foreach \x in {1,...,6} {
            \pgfmathtruncatemacro\r{round(\x-1)}
            \pgfmathtruncatemacro\ethLastEnd{round((\x-1)*\ChainBBlocksN)}
            \pgfmathtruncatemacro\ethBStart{round((\x-1)*\ChainBBlocksN+1)}
            \pgfmathtruncatemacro\ethBEnd{round((\x-1)*\ChainBBlocksN+\ChainBBlocksN)}
            \xdef\stateXS{5pt}

            %%%%%%%%%
            % bitcoin (well, chain A. the diagram grew out of an example of BTC/ETH, you see)
            %%%%%%%%%

            \ifnum 1=\x
                \xdef\ys{0}
                \xdef\prevNode{btcDots}
                \xdef\oPrev{ethDots}
            \else
                % \xdef\ys{-1cm}
                \xdef\prevNode{btc\r}
                \xdef\oPrev{\ethE\r}
            \fi

            \ifnum \DrawBtcLT>\x
                \node (btc\x) [block, below = of \prevNode |- \oPrev, yshift=\ys] {$\AbbrChainA_{i+\x}$};
                \draw [arrowParent] (btc\x) -- (\prevNode);

                % btc state
                \def\InitStateContents{
                    \textsc{STATE} \\
                    Headers: $\AbbrChainB_{0 \cdots j+\ethLastEnd}$
                }
                \ifnum 1=\DrawPoRInState
                    \def\StateContents{\InitStateContents \\
                        Latest PoR: $\AbbrChainA_{i+\r}$ \\
                        %% $\Sigma$ Weight: $\Sigma^{j+\ethLastEnd}_{k=0} \text{W}(\AbbrChainB_{k})$
                    }
                \else
                    \def\StateContents{\InitStateContents}
                \fi

                \ifnum 1=\DrawBtcState
                    \draw [braceLR] (btc\x.west |- btc\x.north) --
                    node (btcState\x) [state, anchor=east, xshift=-\stateXS, align=right] {\footnotesize{\StateContents}}
                    (btc\x.west |- btc\x.south);
                \fi
                % end btc state

                \def\porLabel{}
                \ifnum 1=\DrawPoRLabel
                    \def\porLabel{\scriptsize{and \ethPorIncl}}
                \fi

                % reflection arrow, only on nodes 2 and beyond
                \ifnum 1=\DrawBtcRefl
                    \ifnum 1<\x
                        \draw [arrowRefl] (btc\x) --
                        node [anchor=center, above, sloped] {\scriptsize{include \ethPorHeaders}}
                        node [anchor=center, below, sloped] {\porLabel}
                        (\ethAll\r);
                    \fi
                \fi
            \fi

            %%%%%%%%%%
            % ethereum
            %%%%%%%%%%

            \xdef\ys{0}

            \ifnum 1=\x
                \xdef\prevNode{ethDots}
                \xdef\oPrev{btc1}
            \else
                \xdef\prevNode{\ethE\r}
                \xdef\oPrev{btc\x}
            \fi

            % exclude last Eth chunk
            \ifnum \DrawEthLT>\x
                \node (ethS\x) [block, below = of \prevNode |- \oPrev, yshift=\ys] {$\AbbrChainB_{j+\ethBStart}$};
                \draw [arrowParent] (ethS\x) -- (\prevNode);

                \ifnum 1<\ChainBBlocksN
                    \node (ethDots\x) [dots, below = of ethS\x, yshift=0.5cm] {\vdots};
                    \draw [arrowParent] (ethDots\x) -- (ethS\x);

                    \node (ethE\x) [block, below = of ethDots\x, yshift=0.5cm] {$\AbbrChainB_{j+\ethBEnd}$};
                    \draw [arrowParent] (ethE\x) -- (ethDots\x);
                \fi

                % Ethereum State
                \def\InitStateContents{
                    \textsc{STATE} \\
                    Headers: $\AbbrChainA_{0 \cdots i+\x}$
                }
                \ifnum 1=\DrawPoRInState
                    \def\StateContents{\InitStateContents \\
                        Latest PoR: $\AbbrChainB_{j+\ethLastEnd}$\\
                        %% $\Sigma$ Weight: $\Sigma^{j+\ethLastEnd}_{k=0} \text{W}(\AbbrChainB_{k})$
                    }
                \else
                    \def\StateContents{\InitStateContents}
                \fi
                \ifnum 1=\DrawEthState
                    \draw [braceRL] (ethS\x.east |- ethS\x.north) --
                    node (ethStateCoord\x) [state, anchor=west, xshift=\stateXS] {\footnotesize{\StateContents}}
                    (ethS\x.east |- ethS\x.south);
                    %\node (ethState\x) [block, right = of ethStateCoord\x, xshift=0cm] {State};
                \fi
                % end Ethereum State


                \ifnum 1<\ChainBBlocksN
                    % curly brace for lots of eth blocks
                    \draw [ethBraceL] (ethS\x.west -| \ethE\x.west) -- node (ethAll\x) [black, midway, yshift=14pt] {
                        %\scriptsize{$\AbbrChainB_{40(j+\x) \cdots 40(j+\x)+39}$}
                        \scriptsize{\ChainBBlocksN\:Blocks}
                    } (\ethE\x.west);
                \fi

                \def\porLabel{}
                \ifnum 1=\DrawPoRLabel
                    \def\porLabel{\scriptsize{and PoR}}
                \fi

                % eth refl
                \ifnum 1=\DrawEthRefl
                    \draw [arrowRefl] (ethS\x) --
                    node [anchor=center, above, sloped] {\scriptsize{include header}}
                    node [anchor=center, below, sloped] {\porLabel}
                    (\oPrev);
                \fi
            \fi
        }%

    \ifnum \DrawBtcLT>\DrawEthLT
        \def\lowNode{btc\lastBtc}
    \else
        \def\lowNode{\ethE\lastEth}
    \fi

    \xdef\xs{-0.2cm}
    \def\leftNode{btcChain}
    \ifnum 1=\DrawBtcState
        \xdef\xs{0.8cm}
        \def\leftNode{btcState1}
    \fi

    % \ifnum 1=\DrawBtcState
    %     \node (topLeft) [left = of btcState1.west |- btcChain.north, minimum width=0, xshift=\xs] {};
    %     \node (botLeft) [left = of btcState1.west |- \lowNode.south, minimum width=0, xshift=\xs] {};
    % \else
    %     \node (topLeft) [left = of btcChain.west |- btcChain.north, minimum width=0, xshift=\xs] {};
    %     \node (botLeft) [left = of btcChain.west |- \lowNode.south, minimum width=0, xshift=\xs] {};
    % \fi

    \node (topLeft) [left = of \leftNode.west |- btcChain.north, minimum width=0, xshift=\xs] {};
    \node (botLeft) [left = of \leftNode.west |- \lowNode.south, minimum width=0, xshift=\xs] {};

    \draw [arrowTime] (topLeft) node () [anchor=west, yshift=-10pt] {\textit{Time}} -- (botLeft);
\end{tikzpicture}
\end{document}
